\SourceLeaf{419}%
\noindent so wird
\[
(P\cos\varphi^2+2Q\cos\varphi\sin\varphi+R\sin\varphi^2)ds=(p\sin\varphi-q\cos\varphi)d\varphi,
\]
also
\[
\frac1r=\frac{P\cos\varphi^2+2Q\cos\varphi\sin\varphi+R\sin\varphi^2}{p\sin\varphi-q\cos\varphi}
\]
oder, da
\[
\cos\varphi=\frac{\mp q}{\BareRoot{(pp+qq)}},\qquad\sin\varphi=\frac{\pm p}{\BareRoot{(pp+qq)}}:
\]
\[
\pm\frac1r=\frac{Pqq-2Qpq+Rpp}{(pp+qq)^{\frac32}}.
\]

\begin{center}4.\end{center}
\noindent\hspace*{1.2em}Das zweideutige Zeichen in der letzten Formel könnte anfangs auffallen, findet sich jedoch bei näherer Betrachtung ganz in der Ordnung. In der That, da dieser Ausdruck bloss von den partiellen Differentialien von $V$ abhängt und die Function $V$ an sich bloss die Natur der krummen Linie darstellt, ohne zugleich den Sinn anzugeben, in welchem sie als beschrieben gedacht wird, so muss die Frage, ob die Curve nach der rechten oder linken Seite convex ist, so lange unentschieden bleiben, bis dieser Sinn anders woher bestimmt ist. Eben so verhält es sich mit der Bestimmung von $\varphi$ durch seine Tangente, wobei für sich eine Unbestimmtheit von $180^{0}$ zurückbleibt. Die Angabe, in welchem Sinn die Curve beschrieben wird, kann auf verschiedene Arten gemacht werden.

I. Durch das Zeichen der Veränderung von $x$. Ist $x$ im Zunehmen, so muss $\cos\varphi$ positiv sein; die obern Zeichen gelten daher, wenn $q$ einen negativen, die untern, wenn $q$ einen positiven Werth hat; beim Abnehmen von $x$ ist es umgekehrt.

II. Durch das Zeichen der Veränderung von $y$. Ist $y$ im Zunehmen, so schreibt ein positiver Werth von $p$ die obern, ein negativer die untern vor; umgekehrt bei abnehmendem $y$.

III. Durch das Zeichen des Werths, welchen die Function $V$ erhält, wenn man die Curve verlässt. Es seien die Variationen von $x,y$, wenn man die Curve senkrecht auf die Tangente nach der Rechten, also in der Richtung $\varphi+90^{0}$, verlässt, $\delta x,\delta y$, und die Grösse dieser Normale $=\delta\varrho$; man hat dann
